% TOP_PROBLEM: {"id": 8400017, "title": "O15 — Polynomial roots modulo a prime", "category_id": 3, "category": {"id": 3, "name": "graph_theory", "display_name": "Graph Theory", "description": "Problems involving graphs, networks, and their properties.", "slug": "graph-theory", "order_index": 3, "created_at": "2026-07-31T15:26:25.671Z"}, "status": "open", "problem_number": "AMR-083-0017", "set_id": 13}
% FIRST_POSTED: 2026-10-01
% ENTRY_KIND: statement_only
% UnsolvedMath ID 8400017; source catalogue code AMR-083-0017
% !TeX program = lualatex
% Definition and sources only; this file is not an LLM solution attempt.
\documentclass[11pt,a4paper]{article}
\usepackage[margin=25mm]{geometry}
\usepackage{fontspec}
\setmainfont{lmroman10-regular.otf}[BoldFont=lmroman10-bold.otf,
 ItalicFont=lmroman10-italic.otf,BoldItalicFont=lmroman10-bolditalic.otf]
\usepackage{amsmath,amssymb,mathtools}
\usepackage{unicode-math}
\setmathfont{latinmodern-math.otf}
\AtBeginDocument{\let\setminus\smallsetminus}
\usepackage{xurl,hyperref}
\hypersetup{colorlinks=true,urlcolor=blue}
\setcounter{secnumdepth}{0}
\setlength{\parindent}{0pt}
\setlength{\parskip}{5pt}
\setlength{\emergencystretch}{3em}
\begin{document}
\section*{O15 — Polynomial roots modulo a prime}\label{8400017}
{\small UnsolvedMath ID 8400017 · AMR-083-0017 · Graph Theory · AMR Open Problem Lists}
\subsection{Definitions and mathematical statement}
For a prime $p$, let $\mathbb F_p=\mathbb Z/p\mathbb Z$. The input consists of $p$ in binary and a polynomial
\[
f(X)=a_0+a_1X+\cdots+a_dX^d\in\mathbb F_p[X],
\]
encoded by its full coefficient list with representatives $0\le a_i<p$. The input is promised to have a root in $\mathbb F_p$. A valid output is an integer $u\in\{0,\ldots,p-1\}$ for which $f(u)=0$ in that field.

Can this search be performed by a deterministic classical algorithm in time polynomial in the dense input length, equivalently polynomial in $(d+1)\log(2p)$? Both the prime and the polynomial degree are variable parameters. The zero polynomial may be assigned the answer zero; a nonzero constant polynomial is outside the root-existence promise.

The algorithm need only output one root. It is not required to factor the entire polynomial, enumerate all roots, or recognize when no root exists. Repeated roots are allowed, and a root need not be simple. No factorization or splitting field representation is supplied with the coefficients.

Scanning every element of $\mathbb F_p$ is a finite procedure but can be exponential in $\log p$. Likewise, polynomial time in the numerical value of $p$ is not the requested bound. The statement uses dense coefficient input deliberately: a sparse or circuit encoding would define a different complexity problem when the degree is exponentially large relative to the encoding.
\subsection{Short English statement}
Can one root of a polynomial over a prime field be found deterministically in polynomial time whenever a root exists?
\subsection{Sources}
\begin{itemize}
\item[S1] ULAM.AI and the UnsolvedMath Contributors, supplied problems.json, record 8400017 (AMR-083-0017), AMR Open Problem Lists. Catalogue identity and source transcription; CC BY 4.0.
\url{https://huggingface.co/datasets/ulamai/UnsolvedMath}.
\item[S2] Adleman \& McCurley - Open Problems in Number Theory Complexity II (1994), O15, p12 / LNCS p302. Original source indexed by AMR-083-0017.
\url{https://mccurley.org/papers/open.ps.gz}.
\end{itemize}
\end{document}
